\documentclass[12pt]{ctexart}
\usepackage[a4paper,margin=1in]{geometry}
\usepackage{amsmath,amssymb,amsthm,mathtools,bm}
\usepackage{booktabs,longtable,array,multirow}
\usepackage{enumitem}
\usepackage{setspace}
\usepackage[T1]{fontenc}
\usepackage[utf8]{inputenc}
\usepackage[hidelinks]{hyperref}

\setstretch{1.2}

\title{MAH Model Note for Coauthors\\
Bellman, State Variables, and a Formal English Model Description}
\author{}
\date{}

\begin{document}
\maketitle
\tableofcontents
\newpage

\part*{Part I. 中文说明：控制变量、状态变量与 Bellman 写法}
\addcontentsline{toc}{part}{Part I. 中文说明：控制变量、状态变量与 Bellman 写法}


\section{说明目的}

这份说明回答五个彼此相关的问题：

\begin{enumerate}[leftmargin=2em]
    \item 为什么在 MAH 这篇文章里，劳动力和资本通常不应作为动态 Bellman 的核心控制变量，而应先在静态块中优化掉？
    \item “生产要素是静态的”具体是什么意思？对医药企业而言，这些生产要素包括什么？
    \item 当前 Bellman 形式中各个符号分别代表什么经济对象？
    \item 如果实证上观察到 MAH 之后广义创新都上升，企业能力状态到底应该如何设定：一维还是二维？
    \item 为什么可以把基准状态写成
    \[
    z=\text{research-commercialization capability},
    \]
    这样做的优点、缺点、文献依据和可识别性分别是什么？
\end{enumerate}

全文的核心观点可以概括为一句话：

\begin{quote}
如果目标是 JDE 甚至更高层级期刊，最稳的基准模型应当把 \emph{research intensity}、
\emph{commercialization route choice}、\emph{organization choice} 和
\emph{entry/exit} 作为动态控制变量；把 labor、capital、materials、
manufacturing quantity 视为静态生产投入，先嵌入优化后的流量利润
\(\Pi_s\)；并用一个一维 latent capability \(z\) 来概括医药企业在研发与商业化执行上的综合能力。
\end{quote}

\section{一、为什么静态生产投入通常应先优化掉}

\subsection{1.1 基本逻辑}

动态企业模型通常区分两类选择：

\begin{itemize}[leftmargin=2em]
    \item \textbf{静态选择}：影响当期利润，但不直接改变下一期状态；
    \item \textbf{动态选择}：影响未来状态、未来价值或未来可行行动集合。
\end{itemize}

在你的论文里，真正需要进入 Bellman 递归的对象，是那些会改变未来价值的变量，例如：

\begin{itemize}[leftmargin=2em]
    \item 研发强度（research intensity）；
    \item commercialization 路径选择（externalize / transfer / abandon / integrate）；
    \item 组织状态选择（stay / switch / exit）；
    \item 进入决策与退出决策。
\end{itemize}

相反，劳动力、资本、中间投入、批次生产量等对象，通常首先影响的是当期经营利润。
如果这些对象不带有动态调整成本，也不直接成为下一期状态变量，那么最自然的处理方式是：

\begin{quote}
先在当期利润函数中把它们优化掉，再把优化后的流量收益放进 Bellman 方程。
\end{quote}

\subsection{1.2 什么叫“先优化掉”}

以 \(A\) 型企业为例，假设其最原始的当期利润写成：

\[
\pi_A(z;n,k,m,x_A)
=
p\cdot q_A(n,k,m;z)
-wn-rk-c_m m
+\Omega_A(z)x_A
-\frac{\chi_A}{\eta}x_A^\eta
-f_A.
\]

这里：

\begin{itemize}[leftmargin=2em]
    \item \(n\)：劳动力投入；
    \item \(k\)：资本服务或设备使用；
    \item \(m\)：中间投入；
    \item \(x_A\)：研发强度；
    \item \(z\)：企业能力状态。
\end{itemize}

如果 \(n,k,m\) 只影响当期利润，不进入下一期状态，那么先解：

\[
\bar{\pi}_A(z;w,r,c_m)
=
\max_{n,k,m}
\left\{
p\cdot q_A(n,k,m;z)-wn-rk-c_m m
\right\}.
\]

这一步做完之后，再写：

\[
\Pi_A(z)
=
\bar{\pi}_A(z;w,r,c_m)
+
\max_{x_A\ge 0}
\left\{
\Omega_A(z)x_A-\frac{\chi_A}{\eta}x_A^\eta
\right\}
-f_A.
\]

这就是“静态层先优化掉，嵌入 \(\Pi_A\)”的含义。

\subsection{1.3 为什么这样做更稳}

这样做的三个好处是：

\begin{enumerate}[leftmargin=2em]
    \item 它把模型的重心放在你真正想识别的动态边际上：研发、商业化、组织重构和进入退出；
    \item 它避免把论文误写成一个资本深化或劳动配置模型；
    \item 它与大量动态企业与创新模型的标准写法一致：Bellman 中使用优化后的 flow payoff，而不是在每期递归中反复展开全部静态投入选择。
\end{enumerate}

\subsection{1.4 什么时候不能优化掉}

只有在下列情形下，资本或劳动才值得显式进入动态 Bellman：

\begin{enumerate}[leftmargin=2em]
    \item 它本身就是下一期的状态变量（例如资本存量、产线能力、GMP 资质能力）；
    \item 它存在显著动态调整成本（例如 hiring/firing frictions、学习成本、安装成本）；
    \item 它恰好是论文的主要识别对象（例如研究 MAH 是否改变厂房投资、用工升级或资质投资）。
\end{enumerate}

如果不满足这三条，最稳的做法仍然是静态优化后并入 \(\Pi_s\)。

\section{二、这里的“生产要素”对医药企业具体指什么}

\subsection{2.1 对 \(A\) 型和 \(C\) 型企业}

如果 \(A\) 表示研发与生产一体化企业，\(C\) 表示生产专门化企业，那么“生产要素”最自然包括：

\begin{itemize}[leftmargin=2em]
    \item 制造 labor；
    \item 设备与厂房的资本服务；
    \item API 和其他中间投入；
    \item 批次生产数量；
    \item 质检、放行、仓储与配套运营投入。
\end{itemize}

这些变量首先影响的是当期 operating profit，因此如果它们不构成动态状态，最自然就属于静态块。

\subsection{2.2 对 \(B\) 型企业}

如果 \(B\) 表示研发专门化企业，它并不以内生制造投入为核心，因此它最关键的控制变量通常不是 labor 或 capital，而是：

\begin{itemize}[leftmargin=2em]
    \item 研发强度 \(x_B\)；
    \item 项目筛选或 idea generation 规模；
    \item commercialization 路线选择 \(h\)；
    \item continue / transfer / abandon 的离散决策。
\end{itemize}

所以 \(B\) 型企业更接近“创新与组织路线选择”问题，而不是“工厂生产函数选择”问题。

\section{三、Bellman 写法与每个符号的含义}

\subsection{3.1 静态生产块}

定义：

\[
\bar{\pi}_A(z;w)
=
\max_{n,k,m}
\{\text{operating revenue}-\text{operating cost}\},
\]

\[
\bar{\pi}_C(z;p_m,w)
=
\max_{n,k,m_c}
\{\text{manufacturing revenue}-\text{manufacturing cost}\}.
\]

这里：

\begin{itemize}[leftmargin=2em]
    \item \(\bar{\pi}_A\)：\(A\) 型企业在给定状态 \(z\) 和价格下，静态生产块优化后的利润；
    \item \(\bar{\pi}_C\)：\(C\) 型企业提供受托生产服务时，静态生产块优化后的利润；
    \item \(w\)：工资或综合要素价格；
    \item \(p_m\)：contract-manufacturing service price。
\end{itemize}

\subsection{3.2 \(B\) 型企业的创新/实现块}

定义：

\[
\Pi_B(z)
=
\max_{x_B,h}
\left\{
x_B\Psi_h(z)-\frac{\chi_B}{\eta}x_B^\eta-f_B
\right\}.
\]

这里：

\begin{itemize}[leftmargin=2em]
    \item \(\Pi_B(z)\)：\(B\) 型企业在状态 \(z\) 下的优化后当期流量收益；
    \item \(x_B\)：research intensity / idea-generation intensity；
    \item \(h\)：commercialization route choice；
    \item \(\Psi_h(z)\)：在给定 route \(h\) 下，每个 idea 的期望净收益；
    \item \(\chi_B\)：研发成本参数；
    \item \(\eta>1\)：凸成本曲率参数；
    \item \(f_B\)：固定成本。
\end{itemize}

如果只保留 external commercialization 这条主路线，则可写为：

\[
\Psi_{\text{ext}}(z)
=
\lambda_{\text{ext}}(z)\bigl(R(z)-p_m-\tau_{ext}\bigr).
\]

这里：

\begin{itemize}[leftmargin=2em]
    \item \(\lambda_{\text{ext}}(z)\)：external commercialization 成功率；
    \item \(R(z)\)：成功 commercialize 之后的私人收益；
    \item \(p_m\)：支付给外部生产方的价格；
    \item \(\tau_{ext}\)：external commercialization route-specific governance/compliance cost。
\end{itemize}

\subsection{3.3 动态 Bellman}

定义：

\[
V_s(z)=
\max_{j\in J(s)}
\left\{
\Pi_j(z)-\kappa_{sj}
+\beta E[V_j(z')\mid z,j]
\right\}.
\]

这里：

\begin{itemize}[leftmargin=2em]
    \item \(V_s(z)\)：当前组织状态为 \(s\)，能力状态为 \(z\) 时，从今天开始最优决策下的总贴现价值；
    \item \(J(s)\)：从当前组织状态 \(s\) 出发可行的下一期组织选择集合；
    \item \(j\)：本期选择的组织形式；
    \item \(\Pi_j(z)\)：若选择组织形式 \(j\)，对应的优化后当期流量收益；
    \item \(\kappa_{sj}\)：从 \(s\) 切换到 \(j\) 的组织转换成本；
    \item \(\beta\)：折现因子；
    \item \(E[V_j(z')\mid z,j]\)：在当前 \(z\) 且选择 \(j\) 时，明日 continuation value 的条件期望。
\end{itemize}

一句话解释这条 Bellman：

\begin{quote}
企业今天不是只看“哪个组织形式当前利润最高”，而是看“当前流量收益减去转换成本，再加上未来继续经营的价值”。
\end{quote}

\section{四、动态控制变量到底应该是什么}

\subsection{4.1 基准模型中的控制变量}

如果目标是 JDE 甚至更高，最稳的动态控制变量是：

\begin{itemize}[leftmargin=2em]
    \item \(x_A, x_B\)：研发强度；
    \item \(h\)：commercialization route choice；
    \item \(j\)：organization choice（stay / switch / exit）；
    \item entrant's entry choice：以什么身份进入。
\end{itemize}

也就是说，\textbf{真正该进入 Bellman 的主控制变量应当是 research intensity、commercialization mode 和 organizational choice，而不是 labor 或 capital。}

\subsection{4.2 为什么不是 labor 或 capital}

因为这篇文章真正想识别的是：

\begin{quote}
MAH 是否改变 invention-to-commercialization 的组织路线。
\end{quote}

这不是一篇关于资本深化、就业配置或厂房投资的文章。如果把 labor 和 capital 直接放在动态 Bellman 中，你会把模型中心从“组织路线与创新实现”拖向“生产函数与要素配置”，这会偏离论文主问题。

\section{五、企业能力状态到底应该是什么}

\subsection{5.1 为什么不一定要用 TFP}

在制造业异质企业模型中，常见状态变量是 productivity / efficiency / TFP。
但在医药创新问题里，企业真正 relevant 的异质性往往不是单纯的工厂物理效率，而是：

\begin{itemize}[leftmargin=2em]
    \item 研发效率；
    \item 临床或注册推进能力；
    \item 商业化执行能力；
    \item 组织协调能力；
    \item 监管经验与项目管理经验。
\end{itemize}

因此，医药企业的“能力状态”不必强行写成 manufacturing TFP。

\subsection{5.2 基准模型：一维 latent capability}

最稳的 baseline 是：

\[
z=\text{research-commercialization capability}.
\]

它的含义是：

\begin{quote}
企业把 original-drug idea 变成可商业化产品的综合能力。
\end{quote}

这个对象可以概括：

\begin{itemize}[leftmargin=2em]
    \item idea generation efficiency；
    \item clinical / regulatory execution ability；
    \item project-management ability；
    \item coordinating with external manufacturers；
    \item retaining value from commercialization。
\end{itemize}

\subsection{5.3 这样设定的优点}

\begin{enumerate}[leftmargin=2em]
    \item 它与论文主问题匹配：MAH 主要改的是 innovation-to-commercialization 路线，而不是工厂物理效率；
    \item 它允许一个 parsimonious baseline：用一维 latent state 总结异质性；
    \item 它和大量动态企业文献中“用一维持久状态概括企业异质性”的做法一致。
\end{enumerate}

\subsection{5.4 这样设定的缺点}

\begin{enumerate}[leftmargin=2em]
    \item 它把 research 和 commercialization 混在一起，机制分解不够锋利；
    \item 它与数据的映射不唯一：同样高的 \(z\)，可能是研发强，也可能是协调强；
    \item 如果政策后创新上升，容易混淆“给定 \(z\) 下回报上升”和“\(z\) 本身提升”。
\end{enumerate}

因此，一维 \(z\) 适合作为 baseline，但最好在附录或 robustness 中讨论更细的扩展。

\section{六、如果上二维，最合理的是哪两个维度}

\subsection{6.1 首选方案}

如果数据与计算都足够强，最合理的二维状态是：

\[
z=(z_R,z_C),
\]

其中：

\begin{itemize}[leftmargin=2em]
    \item \(z_R\)：research capability；
    \item \(z_C\)：commercialization / execution capability。
\end{itemize}

为什么这两个维度最合理？因为 MAH 改变的核心不是科学知识生产本身，而是 research-originated ideas 能否被外部实现、推进、协调、注册并最终 commercialize。因此：

\begin{itemize}[leftmargin=2em]
    \item \(z_R\) 管“能否生成高质量 idea”；
    \item \(z_C\) 管“能否把 idea 推进成产品”。
\end{itemize}

\subsection{6.2 备选方案}

另一种可能是：

\[
z=(z_R,z_M),
\]

其中：

\begin{itemize}[leftmargin=2em]
    \item \(z_R\)：research capability；
    \item \(z_M\)：manufacturing / qualified production capability。
\end{itemize}

这个版本更容易和 manufacturing licenses、剂型范围、GMP 资质等数据对接，
但它的缺点是会把论文核心部分地误写成“制造能力问题”，而你的主问题更接近“跨企业 commercialization 的治理与执行能力问题”。

因此，如果必须在两个维度中选最合理的一组，优先建议：

\[
(z_R,z_C).
\]

\subsection{6.3 上二维需要的数据}

至少需要两组 proxy：

\paragraph{对 \(z_R\) 的 proxy}
\begin{itemize}[leftmargin=2em]
    \item R\&D personnel / scientists；
    \item R\&D expenditure；
    \item patent stock / citations；
    \item IND productivity；
    \item novel target share；
    \item 历史 new-drug application success。
\end{itemize}

\paragraph{对 \(z_C\) 的 proxy}
\begin{itemize}[leftmargin=2em]
    \item NDA / approval / launch conversion；
    \item 历史 commercialization success；
    \item regulatory experience；
    \item manufacturing qualification scope；
    \item entrusted-manufacturing relationships；
    \item 持证与外部生产的历史经验；
    \item 质量、放行、合规执行经验。
\end{itemize}

如果这两组数据拿不到，不要急着上二维。

\section{七、关于 \texorpdfstring{$z=\text{research-commercialization capability}$}{z = research-commercialization capability} 的更深讨论}

\subsection{7.1 第一问：概念上可不可以}

可以。因为在动态企业模型里，状态变量的要求不是“必须是直接观测到的 TFP”，而是：

\begin{quote}
它能否作为一个 sufficient statistic，稳定地概括当前企业面对的 payoffs、choices 和 transitions。
\end{quote}

所以，只要 \(z\) 能系统性地影响研发收益、商业化收益、组织切换收益和存活概率，它就可以作为一个合法的 latent capability state。

\subsection{7.2 第二问：顶刊能不能接受}

作为 baseline，完全可以接受，但有前提：

\begin{enumerate}[leftmargin=2em]
    \item 你要明确它是 latent summary state，而不是直接观测的物理 TFP；
    \item 你不能让政策效应本身也偷偷塞进 \(z\) 里；政策仍然应通过独立的 friction 或 expected return shifter 进入；
    \item 你要给出 observable proxies、moments 或 calibration targets 来 discipline 这个状态变量。
\end{enumerate}

\subsection{7.3 第三问：它能不能被数据估计}

它通常不能被直接“看见”，但可以被间接识别或校准。常见方式有三种：

\begin{enumerate}[leftmargin=2em]
    \item \textbf{proxy approach}：用 R\&D personnel、patents、IND productivity、launch conversion 等作为 noisy proxies；
    \item \textbf{moment-based discipline}：让模型匹配 innovation intensity、survival、entry composition、route adoption、conversion 等矩；
    \item \textbf{discrete-state approximation}：把企业按 observable indices 分成 low / medium / high capability bins，再离散化状态。
\end{enumerate}

因此，\(z\) 作为 latent capability 是可估的，但通常是\textbf{间接估计}或\textbf{结构校准}，而不是像会计变量那样直接观测。

\section{八、推荐的基准模型写法}

综合上面的讨论，最稳的基准模型是：

\begin{itemize}[leftmargin=2em]
    \item \textbf{状态变量}：
    \[
    z=\text{research-commercialization capability};
    \]
    \item \textbf{静态生产块}：
    labor、capital、materials、manufacturing quantity 先优化掉，写成 \(\bar{\pi}_A,\bar{\pi}_C\)；
    \item \textbf{动态控制变量}：
    \[
    x_A,\ x_B,\ h,\ j,\ \text{entry/exit};
    \]
    \item \textbf{Bellman}：
    \[
    V_s(z)=\max_{j\in J(s)}
    \left\{
    \Pi_j(z)-\kappa_{sj}+\beta E[V_j(z')\mid z,j]
    \right\}.
    \]
\end{itemize}

如果未来数据与计算都更强，再扩展为：

\[
z=(z_R,z_C).
\]

\section{结论}

这篇 MAH 论文最稳的建模口径不是“制造业 TFP 模型”，而是：

\begin{quote}
一个以 innovation-to-commercialization 组织路线为核心的动态企业模型。
\end{quote}

因此：

\begin{enumerate}[leftmargin=2em]
    \item Bellman 中真正要优化的是 research intensity、commercialization route、organization choice 和 entry/exit；
    \item labor、capital、materials 等生产要素若不带动态调整，就应先在静态利润里优化掉；
    \item 医药企业的基准状态变量不必是 manufacturing TFP，而更合理地是
    \[
    z=\text{research-commercialization capability};
    \]
    \item 如果上二维，首选
    \[
    z=(z_R,z_C).
    \]
\end{enumerate}

这套写法既符合动态企业与创新文献的标准分工，也更贴合 MAH 作为“改变 commercialization route 的制度改革”这一论文主问题。


\newpage
\part*{Part II. English Note: States, Controls, Static Inputs, Prices, and Payoff Objects}
\addcontentsline{toc}{part}{Part II. English Note: States, Controls, Static Inputs, Prices, and Payoff Objects}


\section{Purpose}

This note rewrites the current MAH model in a journal-style format that makes the
economic objects fully explicit. The goal is not to enlarge the model mechanically,
but to clarify exactly which variables are states, which are dynamic controls, which
are static inputs, which are prices, and which objects enter the Bellman equations
as optimized flow payoffs. The note is written to match the conventions of the
dynamic heterogeneous-firm literature and to fit the substantive question of the
paper: how MAH changes the organizational route from invention to commercialization
for original-drug ideas.

\section{Environment}

Time is discrete, indexed by $t=0,1,2,\dots$. There are three organizational types,
\[
s\in\{A,B,C\},
\]
where:

\begin{itemize}[leftmargin=2em]
    \item $A$ denotes integrated firms that can conduct original-drug R\&D and also
    organize production internally;
    \item $B$ denotes research-specialized firms that generate original-drug ideas but
    do not operate as full internal manufacturers;
    \item $C$ denotes manufacturing-specialized firms that provide production capability
    and can supply qualified contract manufacturing.
\end{itemize}

There is also a mass of potential entrants, who may enter as one of the active organizational
types after paying the corresponding entry cost.

Each incumbent firm enters the period with an idiosyncratic capability state
\[
z\in Z\subset \mathbb{R}_+.
\]
In the baseline model, $z$ is one-dimensional. It is interpreted as a broad firm-specific
capability object that loads differently across organizational forms. Economically, $z$
summarizes the firm's ability to generate original-drug ideas, to manage projects,
to coordinate implementation and commercialization, and to adapt organizationally.

\section{Classification of variables}

Table \ref{tab:objects} classifies the main variables and model objects.

\begin{longtable}{p{0.18\textwidth}p{0.20\textwidth}p{0.56\textwidth}}
\caption{Classification of model objects}\label{tab:objects}\\
\toprule
Object & Type & Interpretation \\
\midrule
\endfirsthead
\toprule
Object & Type & Interpretation \\
\midrule
\endhead

$z$ & State variable & Persistent firm capability state. \\[0.4em]

$s\in\{A,B,C\}$ & Discrete state variable & Current organizational form of the firm. \\[0.4em]

$w$ & Price & Wage or composite factor price. \\[0.4em]

$p_m$ & Price & Price of qualified contract-manufacturing services. \\[0.4em]

$n,k,m$ & Static input choices & Labor, capital services, and intermediate inputs used in current-period production. \\[0.4em]

$x_A,x_B$ & Dynamic control variables & Original-drug idea-generation intensities chosen by type-$A$ and type-$B$ firms. \\[0.4em]

$h$ & Dynamic control variable & Commercialization route chosen for a type-$B$ idea, e.g. externalize, transfer, abandon. \\[0.4em]

$j$ & Dynamic control variable & Next organizational choice (stay, switch, exit). \\[0.4em]

$\kappa_{sj}$ & Adjustment-cost object & Organizational switching cost from current type $s$ to chosen type $j$. \\[0.4em]

$\tau_{ext}$ & Route-specific friction & Governance/compliance friction of external commercialization through qualified outside manufacturing. \\[0.4em]

$\bar{\pi}_A(z;w)$ & Optimized static payoff & Current-period operating profit of type $A$ after static production inputs are optimized out. \\[0.4em]

$\bar{\pi}_C(z;p_m,w)$ & Optimized static payoff & Current-period operating profit of type $C$ after static production inputs are optimized out. \\[0.4em]

$\Psi_h(z)$ & Per-idea net value object & Expected net value of one original-drug idea under route $h$. \\[0.4em]

$\Pi_A(z),\Pi_B(z),\Pi_C(z)$ & Optimized flow payoff objects & Current-period payoffs that enter the Bellman equation. \\[0.4em]

$V_A(z),V_B(z),V_C(z)$ & Value functions & Continuation values under each organizational type. \\

\bottomrule
\end{longtable}

\section{Why static production inputs are optimized out}

The model distinguishes between \emph{static} and \emph{dynamic} choices.

\paragraph{Static choices.}
Labor, capital services, and intermediate inputs affect current-period operating profits.
If they do not become next-period states and do not involve dynamic adjustment costs, then
they should not be carried as separate controls in the Bellman equation. Instead, the firm
solves a static profit maximization problem conditional on $(z,s)$ and prices.

\paragraph{Dynamic choices.}
Research intensity, commercialization route, organizational switching, entry, and exit affect
future values, future states, or the feasibility of future actions. These are the natural
dynamic controls.

This separation is standard in dynamic firm models: the Bellman equation is written in
terms of optimized current payoffs rather than raw production problems in every period.
That separation is especially appropriate here because the economic question is about
organizational routes from invention to commercialization, not about capital deepening
or static labor demand per se.

\section{Static production blocks}

\subsection{Integrated firms}

For type-$A$ firms, define the optimized static operating payoff by
\[
\bar{\pi}_A(z;w,r,c_m)
=
\max_{n,k,m}
\left\{
p_A q_A(n,k,m;z)-wn-rk-c_m m
\right\}.
\]
Here:
\begin{itemize}[leftmargin=2em]
    \item $n$ is labor input,
    \item $k$ is capital services or plant use,
    \item $m$ is intermediate input use,
    \item $p_A$ is the relevant output price or revenue shifter for current operations,
    \item $w,r,c_m$ are input prices.
\end{itemize}
This object is \emph{not} the full type-$A$ flow payoff. It is only the optimized
current-period operating block after static production choices have been solved.

\subsection{Manufacturing-specialized firms}

For type-$C$ firms, define the optimized static manufacturing payoff by
\[
\bar{\pi}_C(z;p_m,w,r,c_c)
=
\max_{n,k,m_c}
\left\{
p_m y_C(n,k,m_c;z)-wn-rk-c_c m_c
\right\}.
\]
This is the current-period profit from supplying qualified manufacturing services
after static input choices are optimized out. The price $p_m$ is the contract-manufacturing
service price, so this block directly links the manufacturing side of the industry
to the commercialization problem of type-$B$ firms.

\section{Idea generation and commercialization}

\subsection{Type-$A$ innovation block}

Let $x_A\ge 0$ denote original-drug idea-generation intensity chosen by a type-$A$ firm.
Given state $z$, define the innovation block as
\[
\max_{x_A\ge 0}
\left\{
x_A\Omega_A(z)-\frac{\chi_A}{\eta}x_A^\eta
\right\},
\qquad \eta>1,
\]
where $\Omega_A(z)$ is the per-unit marginal value of original-drug research for
an integrated firm, $\chi_A>0$ is a cost parameter, and $\eta>1$ governs curvature.

Then the optimized type-$A$ flow payoff is
\[
\Pi_A(z)
=
\bar{\pi}_A(z;w,r,c_m)
+
\max_{x_A\ge 0}
\left\{
x_A\Omega_A(z)-\frac{\chi_A}{\eta}x_A^\eta
\right\}
-f_A.
\]

\subsection{Type-$B$ idea-generation and route-choice block}

For type-$B$ firms, the core current-period decision is not static manufacturing.
Instead, it is the combination of idea generation and the route by which ideas
are commercialized.

Let $x_B\ge 0$ denote idea-generation intensity. Let
\[
h\in\mathcal{H}
\]
denote the commercialization route for a type-$B$ idea. In the parsimonious
version of the model, one may let
\[
\mathcal{H}=\{\text{externalize},\text{transfer},\text{abandon}\}.
\]

For any route $h$, let $\Psi_h(z)$ denote the expected net value of one idea.
Then the optimized type-$B$ flow payoff is
\[
\Pi_B(z)
=
\max_{x_B,h}
\left\{
x_B\Psi_h(z)-\frac{\chi_B}{\eta}x_B^\eta-f_B
\right\}.
\]

\subsection{External commercialization as the MAH-sensitive route}

The route that carries the MAH mechanism most directly is external commercialization.
For that route, write
\[
\Psi_{\mathrm{ext}}(z)
=
\lambda_{\mathrm{ext}}(z)\bigl(R(z)-p_m-\tau_{ext}\bigr).
\]

Each term has a distinct interpretation:

\begin{itemize}[leftmargin=2em]
    \item $\lambda_{\mathrm{ext}}(z)$ is the success rate of external commercialization;
    \item $R(z)$ is the gross private return from a successfully commercialized original-drug idea;
    \item $p_m$ is the payment to the qualified external manufacturer;
    \item $\tau_{ext}$ is the route-specific governance/compliance cost of external commercialization.
\end{itemize}

It is essential to distinguish $p_m$ from $\tau_{ext}$. The former is the market price
paid for outside production. The latter is the extra organizational and legal friction of
using that route while retaining the project and authorization. MAH is modeled not as
a generic policy dummy that changes everything at once, but as a reform that lowers
the friction associated with this route.

\section{Dynamic Bellman system}

\subsection{General form}

For each organizational type $s\in\{A,B,C\}$, let the value function be
\[
V_s(z)=
\max_{j\in J(s)}
\left\{
\Pi_j(z)-\kappa_{sj}
+\beta E[V_j(z')\mid z,j]
\right\}.
\]

This equation has four layers:

\begin{enumerate}[leftmargin=2em]
    \item \(\Pi_j(z)\): optimized current-period payoff if the firm operates as type \(j\) this period;
    \item \(\kappa_{sj}\): cost of moving from current organization \(s\) to chosen organization \(j\);
    \item \(\beta\): discount factor;
    \item \(E[V_j(z')\mid z,j]\): expected continuation value next period.
\end{enumerate}

The Bellman equation therefore says that the firm chooses the organizational form
that maximizes the sum of optimized current payoff, minus switching cost, plus discounted
future value.

\subsection{Type-specific Bellman equations}

The general form can be written more explicitly as follows.

\paragraph{Type $A$.}
\[
V_A(z)
=
\max\left\{
\Pi_A(z)+\beta E[V_A(z')\mid z,A],
-\kappa_{AB}+\Pi_B(z)+\beta E[V_B(z')\mid z,B],
-\kappa_{AC}+\Pi_C(z)+\beta E[V_C(z')\mid z,C],
0
\right\}.
\]

\paragraph{Type $B$.}
\[
V_B(z)
=
\max\left\{
\Pi_B(z)+\beta E[V_B(z')\mid z,B],
-\kappa_{BA}+\Pi_A(z)+\beta E[V_A(z')\mid z,A],
-\kappa_{BC}+\Pi_C(z)+\beta E[V_C(z')\mid z,C],
0
\right\}.
\]

\paragraph{Type $C$.}
\[
V_C(z)
=
\max\left\{
\Pi_C(z)+\beta E[V_C(z')\mid z,C],
-\kappa_{CA}+\Pi_A(z)+\beta E[V_A(z')\mid z,A],
-\kappa_{CB}+\Pi_B(z)+\beta E[V_B(z')\mid z,B],
0
\right\}.
\]

\section{Which objects are states, which are controls?}

The model should be read as follows.

\subsection{State variables}

The baseline state variables are:
\[
(s,z).
\]
That is:
\begin{itemize}[leftmargin=2em]
    \item $s$ tells us the current organizational form;
    \item $z$ tells us the current capability state.
\end{itemize}

\subsection{Dynamic controls}

The main dynamic controls are:
\begin{itemize}[leftmargin=2em]
    \item $x_A$ and $x_B$ (research intensities),
    \item $h$ (commercialization route),
    \item $j$ (organization choice),
    \item entry/exit decisions.
\end{itemize}

\subsection{Static inputs}

The main static inputs are:
\begin{itemize}[leftmargin=2em]
    \item labor $n$,
    \item capital services $k$,
    \item intermediates $m$,
    \item manufacturing quantity or manufacturing service intensity where relevant.
\end{itemize}
These inputs are not discarded economically; they are simply optimized out before the
Bellman recursion is written.

\subsection{Prices}

The key equilibrium prices are:
\begin{itemize}[leftmargin=2em]
    \item $w$: labor or composite factor price,
    \item $p_m$: price of qualified contract manufacturing.
\end{itemize}

\subsection{Payoff objects}

The main payoff objects are:
\begin{itemize}[leftmargin=2em]
    \item $\bar{\pi}_A,\bar{\pi}_C$: optimized static profit blocks,
    \item $\Psi_h$: expected per-idea net value under route $h$,
    \item $\Pi_A,\Pi_B,\Pi_C$: optimized current flow payoffs,
    \item $V_A,V_B,V_C$: continuation values.
\end{itemize}

\section{Why the baseline state should be one-dimensional}

The one-dimensional benchmark
\[
z=\text{research-commercialization capability}
\]
is attractive because it is parsimonious and directly tied to the MAH question.
The reform changes the route from invention to commercialization, so the most natural
summary state is not pure manufacturing TFP, but a broader capability object that
shifts research productivity, implementation success, organizational adaptation, and
the value of using external commercialization.

The main gain of this baseline is tractability. It allows the model to capture persistent
heterogeneity without creating multiple latent processes that may not be separately
identified in the data. The main cost is that research capability and commercialization
capability are compressed into a single state. If the data become rich enough, the natural
next step is a two-dimensional state:
\[
z=(z_R,z_C),
\]
where $z_R$ is research capability and $z_C$ is commercialization/execution capability.
But as a baseline intended for estimation or disciplined calibration, the one-dimensional
state is usually the safer choice.

\section{Interpretation for the paper}

The most important implication of this notation is conceptual. The paper is not a model
of capital accumulation, labor adjustment, or static production reallocation. It is a model
of how institutional reform changes the route through which original-drug ideas are turned
into commercialized products. For that reason:

\begin{itemize}[leftmargin=2em]
    \item labor and capital belong to optimized static profit blocks unless they are themselves dynamic states;
    \item research intensity, commercialization route, and organizational switching are the core dynamic controls;
    \item the external commercialization route is the natural place where the MAH friction enters;
    \item the Bellman equations should be written in terms of optimized flow payoffs rather than raw production problems.
\end{itemize}

This is the cleanest way to align the model with the empirical question and with the
conventions of the top dynamic firm and innovation literature.

\end{document}
